% =====================================================================
%  IEEE conference paper - 6-page version (revised content)
%  Compile:  pdflatex "IEEE 6 pages.tex"   (run twice so cross-references settle)
% =====================================================================
\documentclass[conference,a4paper]{IEEEtran}
\IEEEoverridecommandlockouts

\usepackage{cite}
\usepackage{amsmath,amssymb,amsfonts}
\usepackage{graphicx}
\graphicspath{{../}{./}}
\usepackage{textcomp}
\usepackage{booktabs}
\usepackage{array}
\usepackage{multirow}
\usepackage{url}
\usepackage[hidelinks]{hyperref}

\newcommand{\hlt}[1]{\textbf{#1}}

\begin{document}

\title{What Survives a Strict Split? An Empirical Audit of
       Long-Document Screenplay Modelling for IMDb Rating Prediction}

% Blinded for double-blind review.
\author{\IEEEauthorblockN{Anonymous authors}}

\maketitle

% =====================================================================
\begin{abstract}
Most screenplay-to-rating studies report accuracy on a random split of an
archival corpus. We ask how much of that accuracy survives a stricter
split. Our corpus joins $5{,}204$ IMSDb screenplays to IMDb metadata;
after an identity audit we keep $5{,}164$ records and publish the
exclusion ledger. We also find a silent truncation defect: scripts are
cut into $256$-word windows, but the encoder takes at most $254$
WordPiece tokens, so on $1{,}186$ scripts recovered from the project
history $98.4\%$ of windows overflow and $31.9\%$ of tokens are lost.
Keeping representation and learner fixed, we change only the partition:
random $k$-fold, title-family \textsc{GroupKFold}, and a chronological
2020--2025 holdout. The best system, which adds runtime and release year
to the sentence embeddings, scores $R^{2}=0.549$, $0.536$ and $0.309$
under the three protocols; a text-only model usable before a film exists
scores $0.424$, $0.411$ and $0.267$. Release year and decade hold
$11.1\%$ of in-sample split gain yet leave holdout MAE exactly unchanged,
and trained the same way, sentence embeddings do no better than TF-IDF.
The protocol explains more of the spread in reported accuracy than the
model does.
\end{abstract}

\begin{IEEEkeywords}
data leakage, evaluation protocol, sentence embeddings, long-document
NLP, screenplay analysis, temporal generalisation, grouped
cross-validation.
\end{IEEEkeywords}

% =====================================================================
\section{Introduction}

Screenplays are long. The raw scripts in our corpus run to about
$9{,}900$ words on average, while an encoder such as
\texttt{all-MiniLM-L6-v2} reads no more than $256$ tokens at a time. The
usual fix, cheap and common, is to slice each script into overlapping
windows, embed them one by one, and average the vectors. Less attention
goes to whether the evaluation around it yields a number a practitioner
should trust. Corpora built from community archives have two quirks that
a shuffled split quietly rewards.

Duplication is the first. A sequel borrows the vocabulary and structure
of the film before it, and a remake can reuse whole stretches of plot and
dialogue; if one sits in training and the other in test, the test film is
easier than it looks. Time is the second. Older films in the archive
carry higher ratings on average, so a model that knows the release year
can explain a good part of the rating from the calendar alone. Once the
test films come from years the model never saw, which is exactly where a
screenplay-stage predictor operates, that advantage is gone.

Under both sits a smaller, more technical problem. Window length is set
in words, truncation happens in subword tokens, and none of the libraries
involved warns when the token budget overflows. We therefore audit an
existing SBERT-plus-gradient-boosting pipeline rather than propose a new
model, and make three contributions.

\begin{itemize}
\item \hlt{C1: a record-identity protocol and a published exclusion
  ledger.} Identity is checked against four keys (SHA-256 content digest,
  IMDb identifier, title$\times$year, and a normalised title-family
  stem), and every excluded record is accounted for
  (Section~\ref{sec:corpus}).
\item \hlt{C2: a measurement of silent truncation.} We run the encoder's
  own tokeniser over the pipeline's windows for $1{,}186$ recovered
  scripts: $98.4\%$ of the $256$-word windows exceed the $254$-token
  budget, and $31.9\%$ of WordPiece tokens never reach the encoder
  (Section~\ref{sec:tokens}).
\item \hlt{C3: a three-protocol evaluation with paired inference.}
  Representation and learner stay fixed and only the partition changes:
  random $k$-fold, title-family \textsc{GroupKFold}, and a chronological
  2020--2025 holdout. We give bootstrap intervals, paired Wilcoxon tests
  with effect sizes, a block-permutation analysis of which features still
  help on future films, and a like-for-like comparison of sentence
  embeddings against TF-IDF (Sections~\ref{sec:results}
  and~\ref{sec:xai}).
\end{itemize}

We are not claiming a new state of the art. Our main result is the
distance between $R^{2}=0.549$ and $R^{2}=0.309$ on the same corpus,
features and learner, and the further fall to $0.267$ once
post-production features are removed.

% =====================================================================
\section{Related Work}

\hlt{Screenplays and film outcomes.} Scripts were first used as inputs
for forecasting return on investment and box-office
revenue~\cite{eliashberg2014,hunter2016}. Later studies turned to
narrative structure: emotional arcs~\cite{reagan2016,chu2017}, how
characters are portrayed~\cite{ramakrishna2017}, MPAA age suitability and
content severity~\cite{shafaei2020,mohamed2020,zhang2021}, and tags
inferred from plot synopses~\cite{kar2018}. Closest to rating regression
is predicting award nominations from the script~\cite{chiu2020}, which
shares its difficulty: the label is a collective judgement that the text
explains only in part. Outside the script, forecasters have used
production metadata~\cite{sharda2006}, social-media and Wikipedia
activity~\cite{asur2010,oghina2012,mishne2006,mestyan2013},
trailers~\cite{madongo2024}, and the sentiment of reviews or
tweets~\cite{naeem2022,tripathi2023}. Predicting ratings or success from
tabular metadata is a routine applied
exercise~\cite{lee2018,bristi2019,gomes2022,udandarao2024,xie2025}, yet
papers rarely describe their partitions beyond a random seed. Vogel~\cite{vogel2020} points out that theatrical results hinge on
release strategy, competition and marketing, and none of these appears
in a script.

\hlt{Representations for long documents.} Sentence-BERT~\cite{reimers2019}
learns dense sentence embeddings with a Siamese objective; it is the
encoder we audit. Its context window is short, so working at document
level means hierarchical attention~\cite{chalkidis2022}, sparse
attention~\cite{beltagy2020}, or chunking and pooling. Chunk-and-pool
needs no fine-tuning and runs on a CPU, but its pooling is fixed rather
than learned and, as Section~\ref{sec:tokens} shows, its window length is
often given in a unit the encoder does not count in.

\hlt{Leakage and split design.} System rankings can flip under a
different random split, and a split should test the kind of
generalisation a paper claims~\cite{gorman2019,sogaard2021}. Train--test
overlap has been quantified for NLP
benchmarks~\cite{lewis2021,elangovan2021}, a cross-disciplinary survey
lists the main forms of leakage~\cite{kapoor2023}, and temporal
generalisation~\cite{lazaridou2021} and corpus
documentation~\cite{bender2018,dodge2021} each have a literature of their
own. As far as we know, none of this has been applied to a screenplay
corpus.

% =====================================================================
\section{Corpus Integrity and Exclusion Audit}
\label{sec:corpus}

\subsection{Source and identity keys}

The corpus links full screenplays from the Internet Movie Script
Database (IMSDb), which several contributors compiled by hand, to IMDb
metadata. Each of the $5{,}204$ joined records has a title, release year,
decade label, IMDb identifier, runtime, an IMDb rating between $1.0$ and
$10.0$, and a path to the script file. Ratings have mean $5.98$, median
$5.80$ and standard deviation $1.44$. When the archived pipeline ran,
every text was on disk; the current snapshot ships only $356$ of them.
Searching the repository history brought back $1{,}186$ files: $1{,}177$
joined records and nine files the metadata never references. They come
from one contiguous block of the rating-ordered file list, so although
they span release years 1958--2025, their ratings all lie between $5.0$
and $5.7$. Any analysis that needs raw text is limited to this subset.

We check four identity keys, from strictest to loosest. The
\emph{content digest}
$\mathrm{SHA256}(b)$ of a file with bytes $b$ catches the same script
re-uploaded under a new name, which no metadata key would notice; when
two files collide we keep whichever loaded first. The \emph{IMDb
identifier} is the canonical film key. \emph{Title$\times$year} is a
fallback for records whose identifier is missing or malformed. The
\emph{title family} is the loosest key and the one that matters for
splitting. To build it we case-fold the title $t$, strip
non-alphanumeric characters, drop a leading article, and keep deleting
trailing sequel markers (an Arabic numeral, a Roman numeral up to
$\mathrm{X}$, or \texttt{part}~$n$) until nothing changes. \emph{Toy
Story} and \emph{Toy Story 2} both end up in the family \texttt{toy
story}, the grouping variable of Section~\ref{sec:splits}.

\subsection{The exclusion flow}
\label{sec:flow}

Let $\mathcal{R}_0$ be the joined record set, $|\mathcal{R}_0| = 5{,}204$.
The pipeline is a composition of four filters applied in a fixed order,
each defined by a predicate on a record $r$:
\begin{align}
\pi_{\text{len}}(r)  &\equiv |\mathrm{text}(r)| \ge 1000 \ \text{characters},\\
\pi_{\text{hash}}(r) &\equiv \mathrm{SHA256}\bigl(\mathrm{bytes}(r)\bigr)
   \notin \mathcal{H}_{<r},\\
\pi_{\text{year}}(r) &\equiv \mathrm{yr}(r) \in [1900,\, 2025],\\
\pi_{\text{id}}(r)   &\equiv \text{positional identity of } r
   \text{ is unique},
\end{align}
where $\mathcal{H}_{<r}$ is the set of digests already seen in load order.
The surviving sets are nested,
$\mathcal{R}_{k} = \{\, r \in \mathcal{R}_{k-1} : \pi_k(r) \,\}$, and the
realised cardinality chain on this corpus is
\begin{equation}
\label{eq:chain}
5{,}204
\;\xrightarrow{\ \pi_{\text{len}}\ } 5{,}195
\;\xrightarrow{\ \pi_{\text{hash}}\ } 5{,}195
\;\xrightarrow{\ \pi_{\text{year}}\ } 5{,}191
\;\xrightarrow{\ \pi_{\text{id}}\ } 5{,}164 .
\end{equation}

On the $1{,}186$ recovered files the digest finds two duplicate groups;
in both, the second copy is one of the nine unreferenced files, so no
joined record is lost. The check covers $1{,}177$ of the $5{,}204$
records ($22.6\%$). The metadata keys can be computed for every
record; they bound the remaining duplication but cannot rule it out. The
length filter $\pi_{\text{len}}$ was applied when the archived pipeline
ran, while every text was still available.

\subsection{Audit results}

\begin{table}[t]
\caption{Corpus integrity audit and exclusion ledger. Percentages are of
the $5{,}204$ joined records.}
\label{tab:audit}
\centering
\small
\begin{tabular}{lrr}
\toprule
\textbf{Check / exclusion} & \textbf{Count} & \textbf{\%} \\
\midrule
Joined records (IMSDb $\bowtie$ IMDb)      & 5{,}204 & 100.00 \\
Distinct IMDb identifiers                  & 5{,}204 & 100.00 \\
Duplicate title$\times$year pairs          &       0 &   0.00 \\
Title collisions (same title, other year)  &  80 / 165 rec. & 3.17 \\
SHA-256 duplicate digests (1{,}177 rec.)   &       0 & --- \\
\addlinespace[2pt]
Script text $<1{,}000$ characters          &       9 &   0.17 \\
Release year outside $[1900, 2025]$        &       4 &   0.08 \\
Unresolvable positional identity           &      27 &   0.52 \\
\midrule
\textbf{Analysis corpus}                   & \textbf{5{,}164} & \textbf{99.23} \\
Title families / multi-member              & 4{,}979 / 164 & --- \\
\bottomrule
\end{tabular}
\end{table}

Table~\ref{tab:audit} gives the complete ledger for the flow in
Section~\ref{sec:flow}. The $5{,}204$ identifiers are all distinct and no
title$\times$year pair occurs twice, so the join itself is clean and
leakage here does not come from plain duplicate records. Eighty titles
occur more than once ($165$ records), but each collision involves
different identifiers and different years: these are remakes, or
unrelated films that share a title, not copies, so they stay in the
corpus as one family. Four records list release years of $1080$ or
$1088$, obvious typos for $1980$ and $1988$; we remove them instead of
guessing a repair. The derived decade column is in worse shape. The
shipped label encoder knows $14$ decade strings, \texttt{1080s} and
\texttt{1996s} among them, and turns them into an \emph{ordinal}
integer.

We exclude $27$ more records for reproducibility reasons. The archived
embedding matrix stores rows in loader order with no record key, so row
$i$ of the matrix need not be record $i$ of the spreadsheet. We rebuilt
the offset function from $1{,}559$ anchors in the committed split
bookkeeping. Inside the gaps of that reconstruction, $27$ positions fit
more than one assignment, and we drop them rather than guess.

% =====================================================================
\section{Methodology and Representations}
\label{sec:method}

\subsection{Structural screenplay statistics}

Sixteen statistics are computed from the raw, uncleaned script text: six
size measures (\texttt{char\_count}, \texttt{word\_count},
\texttt{line\_count}, \texttt{sentence\_count}, \texttt{scene\_count},
\texttt{words\_per\_scene}); three describing lexical sophistication
(\texttt{avg\_word\_length}, \texttt{unique\_word\_ratio},
\texttt{long\_word\_ratio}); two describing sentence geometry
(\texttt{avg\_sentence\_length}, \texttt{sentence\_length\_std}); two
describing dialogue (\texttt{dialogue\_density}, the share of lines that
match an ALL-CAPS speaker-cue pattern, and \texttt{unique\_characters});
two emotional proxies (\texttt{exclamation\_ratio},
\texttt{question\_ratio}); and one stage-direction proxy
(\texttt{action\_density}). Two of the sixteen are close to degenerate on
this corpus. Mean \texttt{dialogue\_density} is $0.0087$ with standard
deviation $0.0361$, and mean \texttt{unique\_characters} is $3.41$. A
feature film has far more than three speaking parts. The speaker-cue
regular expression assumes a \texttt{NAME:} convention that most IMSDb
transcripts do not follow, so both features encode transcription format
rather than dramaturgy.

\subsection{Token-aligned chunking and mean pooling}
\label{sec:tokens}

Let a screenplay be a word sequence $D = (w_1, \dots, w_N)$ obtained by
whitespace splitting after light normalisation that removes bracketed
stage directions, ALL-CAPS speaker cues, \texttt{INT.}/\texttt{EXT.}
markers, and timestamps. With window length $L = 256$ and overlap
$o = 50$, the stride is $s = L - o = 206$, and the $j$-th window is
\begin{equation}
C_j \;=\; \bigl(w_{(j-1)s+1}, \dots, w_{\min\{(j-1)s+L,\; N\}}\bigr),
\end{equation}
for $j = 1, \dots, M$ with $M = 1$ if $N \le L$ and
$M = \lceil (N - L)/s \rceil + 1$ otherwise. Each window is encoded by
$f: \mathcal{C} \to \mathbb{R}^{d}$ with $d = 384$, and the document
representation is the arithmetic mean
\begin{equation}
\label{eq:meanpool}
e(D) \;=\; \frac{1}{M} \sum_{j=1}^{M} f(C_j) .
\end{equation}
Because the encoder L2-normalises its output, every $f(C_j)$ lies on the
unit sphere, so
\begin{equation}
\label{eq:norm}
\bigl\| e(D) \bigr\|_2^2
= \frac{1}{M^2}\Bigl( M + 2\!\!\sum_{1 \le i < j \le M}\!\!
  \langle f(C_i), f(C_j)\rangle \Bigr),
\end{equation}
which equals $1$ only when all windows embed identically and falls toward
$1/\sqrt{M}$ as they decorrelate. With a median of $33$ windows per
recovered script, any single scene has little influence on the pooled
vector.

\hlt{The word/token mismatch.} The encoder's transformer accepts
$T_{\max} = 256$ positions, two of which are consumed by \texttt{[CLS]}
and \texttt{[SEP]}. The content budget is therefore $B = T_{\max} - 2 =
254$ WordPiece tokens, while $C_j$ is defined by a count of $256$
\emph{words}. Writing $\tau(\cdot)$ for the tokeniser, the encoder in
fact sees $\tau(C_j)$ truncated to its first $B$ entries, and the
discarded fraction for a document is
\begin{equation}
\rho(D) \;=\; 1 - \frac{\sum_{j=1}^{M} \min\bigl\{|\tau(C_j)|,\, B\bigr\}}
                        {\sum_{j=1}^{M} |\tau(C_j)|} .
\end{equation}

Running the encoder's own tokeniser over the pipeline's windows on the
$1{,}186$ recovered scripts yields $44{,}899$ windows with median length
$373$ tokens, $90$th percentile $409$ and maximum $706$. Of these,
$98.4\%$ exceed $B$, and $\rho$ computed over all tokens is $0.319$
(worst script $0.611$). Roughly one word in three is written to a window
position that the encoder never reads, and the pipeline reports nothing.
The $356$-file public slice gives almost the same figures ($98.45\%$ and
$0.316$), so the rate is stable across subsets, although the recovered
files cover only ratings $5.0$--$5.7$.

\subsection{Lexical control and learners}

The lexical control is TF-IDF over unigrams and bigrams, capped at
$8{,}000$ features, with English stop words removed, minimum document
frequency $3$, maximum document frequency $0.85$, and sublinear term
frequency $1 + \log \mathrm{tf}$. It uses the same gradient-boosted
regressor as the semantic arm. In the archived runs, however, the TF-IDF
arm received text alone, whereas the SBERT arm also received the sixteen
structural statistics and the three metadata columns and was trained
with inverse-frequency sample weights. Section~\ref{sec:lexical} reports
a like-for-like rerun.

The regressor throughout is gradient-boosted trees with $500$ rounds,
learning rate $0.05$, maximum depth $6$, histogram splitting, and early
stopping after $20$ rounds without validation improvement; predictions
are clipped to $[1, 10]$. Two linear probes provide reference points:
ordinary least squares on the three metadata columns, and ridge
regression ($\alpha = 1.5$) on standardised embeddings. A constant-mean
predictor anchors the bottom. For the stacked ensemble, a ridge
meta-regressor is fitted over four base models --- metadata OLS,
structural OLS (sixteen statistics plus metadata), TF-IDF+XGBoost,
SBERT+XGBoost --- using inner five-fold out-of-fold meta-features nested
within each outer fold. In Tables~\ref{tab:grouped}--\ref{tab:paired},
SBERT + XGBoost uses the $384$ embedding dimensions only, and SBERT +
meta + XGBoost adds runtime, release year and the decade code ($387$
inputs).

\subsection{Partition protocols}
\label{sec:splits}

Let $\mathcal{I} = \{1, \dots, n\}$ index the analysis corpus,
$n = 5{,}164$.

\hlt{P1, random.} A shuffled $K$-fold partition, $K = 5$: disjoint
$F_1, \dots, F_K$ with $\bigcup_k F_k = \mathcal{I}$, assignment
independent of every record attribute. Within each training pool, a
random $18.75\%$ slice is held out for early stopping; the same applies
to P2.

\hlt{P2, grouped.} Let $\gamma: \mathcal{I} \to \mathcal{G}$ map each
record to its title family, $|\mathcal{G}| = 4{,}979$.
\textsc{GroupKFold} constructs disjoint folds subject to the group
constraint
\begin{equation}
\label{eq:group}
\forall g \in \mathcal{G}\;\; \exists!\, k \in \{1,\dots,K\}:\;
\gamma^{-1}(g) \subseteq F_k ,
\end{equation}
so no family is split across the boundary. Because $4{,}979$ families
cover $5{,}164$ records, the constraint is nearly free: realised fold
sizes are $1{,}032$--$1{,}033$, identical to P1.

\hlt{P3, chronological.} With $\mathrm{yr}(i)$ the release year, the
test set is $2020 \le \mathrm{yr}(i) \le 2025$, validation
$c \le \mathrm{yr}(i) \le 2019$, and training $\mathrm{yr}(i) < c$, where
$c$ is the $85$th percentile of the pre-2020 year distribution,
here $c = 2016$. Early stopping is thereby decided on films that are
themselves in the future relative to the training set, so no
forward-looking information enters through model selection. The realised
sizes are $3{,}918$ / $731$ / $515$. All models, including the linear
baselines, are fitted on the training years only.

A partition's leakage can be measured directly. Define the family-overlap
rate of a test fold as the share of its records whose family also occurs
in the corresponding training pool,
\begin{equation}
\lambda(F_k) = \frac{1}{|F_k|}
\bigl|\{\, i \in F_k : \gamma(i) \in \gamma(\mathcal{I} \setminus F_k) \,\}\bigr| .
\end{equation}
Averaged over folds, $\lambda = 5.50\%$ under P1 and, by construction,
$\lambda = 0$ under P2. The chronological holdout sits at
$\lambda = 6.41\%$ --- a sequel released in 2021 may well have a parent
film in the training years.

\subsection{Inference}

Point estimates carry percentile bootstrap intervals from $B = 1000$
resamples of the test set. Model comparison is paired and
distribution-free. For predictions $\hat{y}^{A}$, $\hat{y}^{B}$ on a
common test set, form per-record absolute errors
$e^{A}_i = |\hat{y}^{A}_i - y_i|$ and $e^{B}_i$ likewise, take
$d_i = e^{A}_i - e^{B}_i$, discard the $n_0$ exact ties, rank the
remaining $|d_i|$, and accumulate
\begin{equation}
R^{+} = \!\!\sum_{i:\, d_i > 0}\!\! \mathrm{rank}(|d_i|),
\qquad
R^{-} = \!\!\sum_{i:\, d_i < 0}\!\! \mathrm{rank}(|d_i|) ,
\end{equation}
with test statistic $W = \min(R^{+}, R^{-})$. The signed-rank test is
non-parametric and has no degrees-of-freedom parameter; the reporting
quantities are the number of pairs $n$, the number of non-tied pairs
$n - n_0$, the statistic $W$, and the matched-pairs rank-biserial effect
size
\begin{equation}
r \;=\; \frac{R^{+} - R^{-}}{R^{+} + R^{-}} \;\in\; [-1, 1] ,
\end{equation}
which we report alongside $p$ so that a small effect at large $n$ is not
mistaken for a large one. Differences in mean absolute error carry their
own paired bootstrap intervals.

% =====================================================================
\section{Experimental Results}
\label{sec:results}

\subsection{The cost of a stricter partition}

Table~\ref{tab:grouped} shows the cost of the group constraint and
Table~\ref{tab:temporal} the cost of the temporal one. The five systems,
the corpus and the learner settings are the same throughout.

\begin{table*}[t]
\caption{Grouped five-fold results. P1 shuffles records; P2 enforces the
title-family constraint \eqref{eq:group}. Pooled out-of-fold point
estimates over $5{,}164$ records. $\Delta R^{2}$ is the explained variance
lost to the constraint, i.e.\ the share of P1 skill that franchise
overlap was supplying.}
\label{tab:grouped}
\centering
\small
\begin{tabular}{lccc@{\hskip 18pt}ccc@{\hskip 18pt}c}
\toprule
& \multicolumn{3}{c}{\textbf{P1 random $k$-fold} ($\lambda = 5.50\%$)}
& \multicolumn{3}{c}{\textbf{P2 grouped by family} ($\lambda = 0\%$)} & \\
\cmidrule(lr){2-4}\cmidrule(lr){5-7}
\textbf{System} & RMSE & MAE & $R^{2}$ & RMSE & MAE & $R^{2}$ & $\Delta R^{2}$ \\
\midrule
Constant mean          & 1.4373 & 1.1761 & $-0.0000$ & 1.4375 & 1.1765 & $-0.0004$ & $-0.0004$ \\
Metadata OLS           & 1.1237 & 0.8701 & 0.3888 & 1.1236 & 0.8695 & 0.3889 & $+0.0001$ \\
SBERT ridge            & 1.1122 & 0.8618 & 0.4012 & 1.1238 & 0.8742 & 0.3887 & $-0.0125$ \\
SBERT + XGBoost        & 1.0907 & 0.8355 & 0.4241 & 1.1026 & 0.8445 & 0.4114 & $-0.0127$ \\
SBERT + meta + XGBoost & \textbf{0.9652} & \textbf{0.7276} & \textbf{0.5490}
                       & \textbf{0.9791} & \textbf{0.7309} & \textbf{0.5359} & $-0.0131$ \\
\addlinespace[2pt]
\multicolumn{8}{l}{\footnotesize Best system: P1 $R^{2}\in[0.5298,0.5671]$,
P2 $R^{2}\in[0.5162,0.5537]$. Every arm that reads text loses $R^{2}$;
the metadata-only probe does not.}\\
\bottomrule
\end{tabular}
\end{table*}

\begin{table}[t]
\caption{Temporal holdout. Fit on films released up to $2015$, early
stopping on $2016$--$2019$, tested on the $515$ films released
$2020$--$2025$.}
\label{tab:temporal}
\centering
\small
\begin{tabular}{lccc}
\toprule
\textbf{System} & RMSE & MAE & $R^{2}$ \\
\midrule
Constant mean   & 1.2230 & 0.9762 & $-0.0947$ \\
Metadata OLS    & 1.0794 & 0.8384 & 0.1473 \\
SBERT ridge     & 1.0719 & 0.8417 & 0.1591 \\
SBERT + XGBoost & 1.0007 & 0.7469 & 0.2671 \\
SBERT + meta + XGBoost & \textbf{0.9717} & \textbf{0.7232} & \textbf{0.3089} \\
\addlinespace[2pt]
\multicolumn{4}{l}{\footnotesize Best system $R^{2}\in[0.1985,0.4038]$;
$\lambda = 6.41\%$.}\\
\multicolumn{4}{l}{\footnotesize Target mean $\pm$ SD: train ($\le 2015$)
$6.11 \pm 1.45$; pre-2020}\\
\multicolumn{4}{l}{\footnotesize pool $6.01 \pm 1.46$; holdout
$5.75 \pm 1.17$.}\\
\bottomrule
\end{tabular}
\end{table}

Grouping costs little. The best system moves from $R^{2} = 0.5490$ to
$0.5359$, a drop of $0.0131$, and the two bootstrap intervals overlap
substantially. This is expected for a corpus in which only $164$
families have more than one member and the largest has four. The
direction of the change is nonetheless consistent: every system that
reads text loses $R^{2}$ (SBERT ridge by $0.0125$, SBERT + XGBoost by
$0.0127$), while the metadata-only probe does not move ($0.3888$ to
$0.3889$). Franchise leakage therefore enters through the text, which a
metadata model never sees.

Chronology costs much more. The best system falls to $R^{2} = 0.3089$,
and its interval $[0.1985, 0.4038]$ no longer overlaps either in-corpus
estimate. The constant-mean predictor drops from $R^{2} \approx 0$ to
$-0.0947$, a pure distribution-shift effect: the training mean ($6.11$,
films up to 2015) is too high for the holdout ($5.75$). The metadata
probe degrades most, from $0.3889$ to $0.1473$, losing $62\%$ of its
explained variance, while the text-only SBERT + XGBoost model loses
$35\%$ ($0.4114$ to $0.2671$). That text-only figure, $0.267$, is also
the most relevant number for screenplay-stage use, because it relies on
nothing that becomes known only after production.

\subsection{Paired comparisons}

Table~\ref{tab:paired} reports every baseline against the best system on
matched records.

\begin{table}[t]
\caption{Paired comparison against SBERT + metadata + XGBoost.
$\Delta\mathrm{MAE}$ is positive when the reference system has lower
error. $p$ from the two-sided Wilcoxon signed-rank test.}
\label{tab:paired}
\centering
\small
\setlength{\tabcolsep}{3pt}
\begin{tabular}{llrrr}
\toprule
\textbf{Prot.} & \textbf{Baseline} & $\Delta\mathrm{MAE}$ & \textbf{95\% CI} & $p$ \\
\midrule
\multirow{4}{*}{P1}
 & Constant mean   & $+0.4485$ & $[+0.427, +0.471]$ & $2.2\!\times\!10^{-281}$ \\
 & Metadata OLS    & $+0.1425$ & $[+0.129, +0.157]$ & $1.8\!\times\!10^{-73}$ \\
 & SBERT ridge     & $+0.1342$ & $[+0.116, +0.150]$ & $7.0\!\times\!10^{-55}$ \\
 & SBERT + XGB     & $+0.1079$ & $[+0.093, +0.122]$ & $4.7\!\times\!10^{-46}$ \\
\addlinespace[2pt]
\multirow{4}{*}{P2}
 & Constant mean   & $+0.4456$ & $[+0.424, +0.466]$ & $2.9\!\times\!10^{-280}$ \\
 & Metadata OLS    & $+0.1386$ & $[+0.124, +0.152]$ & $2.8\!\times\!10^{-73}$ \\
 & SBERT ridge     & $+0.1433$ & $[+0.127, +0.159]$ & $2.5\!\times\!10^{-61}$ \\
 & SBERT + XGB     & $+0.1136$ & $[+0.100, +0.127]$ & $2.1\!\times\!10^{-50}$ \\
\addlinespace[2pt]
\multirow{4}{*}{P3}
 & Constant mean   & $+0.2530$ & $[+0.195, +0.313]$ & $5.7\!\times\!10^{-15}$ \\
 & Metadata OLS    & $+0.1152$ & $[+0.063, +0.166]$ & $2.0\!\times\!10^{-4}$ \\
 & SBERT ridge     & $+0.1185$ & $[+0.064, +0.176]$ & $1.3\!\times\!10^{-4}$ \\
 & SBERT + XGB     & $+0.0236$ & $[-0.020, +0.074]$ & $0.273$ \\
\bottomrule
\end{tabular}
\end{table}

The last row matters most. Under P1 and P2, adding the three metadata
columns to the embedding lowers MAE by $0.108$--$0.114$ rating points,
with very small $p$-values. Under P3 the same addition lowers MAE by only
$0.0236$, the interval includes zero, and $p = 0.273$ ($W = 62{,}734$
over $n = 515$ pairs, no ties, rank-biserial $r = +0.056$). By the usual
conventions that effect is negligible. Metadata that looks highly
significant under random folds brings no measurable gain on future
films.

\subsection{The lexical and structural arms}
\label{sec:lexical}

The lexical arm and the structural model need raw text for the whole
corpus, so they exist only as archived runs on the original random
splits over $5{,}195$ records: five folds, plus one $70/15/15$ split
whose test predictions were saved. The archived structural OLS uses the
sixteen statistics together with the three metadata columns; it reaches
$R^{2} = 0.440 \pm 0.007$ against $0.382$ for metadata alone, so the
statistics add about $0.06$. The archived SBERT arm also received all
nineteen numeric columns, whereas TF-IDF + XGBoost saw text only, so its
margin over TF-IDF mixed representation with metadata. For a fair
comparison we retrained an SBERT-only model on
the archived embeddings with the TF-IDF arm's exact recipe ($300$ trees,
no early stopping, the same sample weights, folds and seeds). Over the
five folds, SBERT-only reaches $R^{2} = 0.372 \pm 0.012$ against
$0.433 \pm 0.014$ for TF-IDF (MAE $0.880$ against $0.819$). On the saved
split, where errors can be paired over $780$ test films, TF-IDF has
lower MAE by $0.030$ $[-0.010, 0.069]$ ($p = 0.088$). Mean-pooled
sentence embeddings therefore do not outperform a bag-of-words model
here; the archived advantage came from the metadata. In the original
pipeline, inverse-frequency sample weighting hurts: removing it lowers
MAE by $0.0482$ $[0.039, 0.057]$
($p = 2.5\times10^{-26}$). A nested ridge stack reaches
$R^{2} = 0.577 \pm 0.010$, improving on its strongest base by
$\Delta\mathrm{MAE} = 0.0134$ $[0.008, 0.019]$, and a $100$-trial TPE
search reaches $R^{2} = 0.581$ $[0.547, 0.618]$ on a single held-out
split. Both gains are small next to the $0.240$ gap in $R^{2}$ between
the random and chronological protocols.

\subsection{Temporal shift in detail}

The holdout target is compressed relative to the training years:
standard deviation $1.17$ against $1.45$, mean $5.75$ against $6.11$.
Part of the fall in $R^{2}$ is therefore arithmetic, since the same
absolute error against a narrower target yields a lower $R^{2}$. MAE
alone does not settle the matter: the best system's MAE is $0.7276$
under P1 and $0.7232$ under P3, but the constant predictor's MAE also
falls, from $1.176$ to $0.976$. Relative to that baseline, the best
system removes $38\%$ of baseline MAE under P1 and only $26\%$ under P3,
so accuracy does degrade on future films. Regressing prediction on truth
over the holdout gives slope $0.463$ (s.e.\ $0.028$): the model separates
above-average from below-average films but underestimates how far they
lie from the mean.

% =====================================================================
\section{What the Features Are Doing}
\label{sec:xai}

\subsection{In-sample gain is not out-of-sample use}

Split gain here comes from the archived SBERT model that saw all
nineteen numeric inputs ($403$ features in total, random split).
Embedding dimensions account for $55.5\%$ of the total gain, the metadata
columns for $36.6\%$ (\texttt{movie\_length} $25.5\%$, \texttt{year}
$11.0\%$, \texttt{decade\_encoded} $0.06\%$) and the structural
statistics for $8.0\%$. Looked at per input, the order flips. The
strongest embedding dimension contributes $1.68\%$ and a typical one
$0.14\%$, while runtime on its own is worth three times as much as all
sixteen structural statistics put together.

\subsection{Block permutation on the temporal holdout}

Gain is measured on training splits, so it keeps every in-sample
advantage; permutation on the holdout does not. For the P3 SBERT + meta
+ XGBoost model we shuffled each feature block across the $515$ holdout
films, re-predicted, and recorded the rise in MAE over $20$ repetitions
(baseline $\mathrm{MAE} = 0.7232$):
\begin{itemize}
\item permuting all three metadata columns: $+0.2168 \pm 0.0250$;
\item permuting \texttt{movie\_length} alone: $+0.1988 \pm 0.0185$;
\item permuting the entire $384$-dimensional embedding:
      $+0.1950 \pm 0.0217$;
\item permuting \texttt{year}: $+0.0000 \pm 0.0000$;
\item permuting \texttt{decade\_encoded}: $+0.0000 \pm 0.0000$.
\end{itemize}

Neither zero is a rounding artefact. All holdout films came out in
2020--2025 and the model only saw films up to 2015, so every holdout
value of \texttt{year} falls past the largest split threshold in the
training data. The trees send all of them down the same branch, and
shuffling values that already land in the same leaf has no effect. A
feature that carries $11\%$ of in-sample gain is, at deployment, a
constant.

Runtime behaves differently. A $130$-minute film from 2023 sits in the
same part of feature space as a $130$-minute film from 1994, so runtime
carries over across time, and by itself it is worth as much as the whole
embedding block. But runtime belongs to the finished film and is unknown
while a screenplay is assessed, so it cannot help screenplay-stage
prediction.

% =====================================================================
\section{Threats to Validity and Limitations}
\label{sec:limits}

Raw text survives for just $1{,}177$ of the $5{,}204$ records, all rated
$5.0$--$5.7$, so the hash audit and the truncation estimate rest on a
narrow rating range, and TF-IDF and the structural statistics exist only
as archived random-split runs. The
embedding matrix cannot be rebuilt from the current snapshot either.
IMSDb's selection bias cannot be corrected after the fact: the mean
rating drops from $7.87$ for 1920s films to $5.48$ for the 2010s, which
says more about what people chose to upload than about the films. We
evaluate on a single holdout window; rolling-origin evaluation would show
how quickly accuracy decays. Title families miss remakes and spin-offs
released under a different title, so P2 removes only part of the
franchise overlap. Finally, we audit only mean-pooled
\texttt{all-MiniLM-L6-v2}, and other encoders may truncate differently.

% =====================================================================
\section{Conclusion}

We set out to audit a screenplay-to-rating pipeline, not to improve it.
At the record level the corpus is clean, but it holds $164$
multi-member title families, four corrupt release years and a broken
decade vocabulary that passes into the feature matrix untouched. The
chunking setup silently throws away about $32\%$ of WordPiece tokens, and
when trained the same way the sentence embeddings do not beat TF-IDF.
Above all, the headline metric depends more on the partition than on the
model: $R^{2} = 0.549$ with random folds, $0.536$ with franchises kept
together, $0.309$ on unseen years and $0.267$ from text alone. We
recommend grouping by title family, holding out by release year,
reporting baseline-relative error under target shift, aligning chunks
with the tokeniser, and keeping an immutable record key next to every
cached representation.

% =====================================================================
\section*{Reproducibility}

Code, the archived embeddings and all result files are at
\url{https://github.com/efazulkarim/imdb-predictor/tree/sbert}; every
number regenerates from the committed JSON artifacts (seed $42$). The
screenplays are not redistributed, but IMDb identifiers and SHA-256
digests are enough to rebuild the corpus.

% =====================================================================
\begin{thebibliography}{99}
\footnotesize

\bibitem{eliashberg2014}
J.~Eliashberg, S.~K.~Hui, and Z.~J.~Zhang, ``Assessing box office
performance using movie scripts: A kernel-based approach,''
\emph{IEEE Trans.\ Knowl.\ Data Eng.}, vol.~26, no.~11, pp.~2639--2648,
2014.

\bibitem{hunter2016}
S.~D.~Hunter, S.~Smith, and S.~Singh, ``Predicting box office from the
screenplay: A text analytical approach,'' \emph{J.\ Screenwriting},
vol.~7, no.~2, 2016.

\bibitem{reagan2016}
A.~J.~Reagan, L.~Mitchell, D.~Kiley, C.~M.~Danforth, and P.~S.~Dodds,
``The emotional arcs of stories are dominated by six basic shapes,''
\emph{EPJ Data Science}, vol.~5, no.~31, 2016.

\bibitem{chu2017}
E.~Chu and D.~Roy, ``Audio-visual sentiment analysis for learning
emotional arcs in movies,'' in \emph{Proc.\ IEEE ICDM}, 2017.

\bibitem{ramakrishna2017}
A.~Ramakrishna, V.~R.~Mart{\'i}nez, N.~Malandrakis, K.~Singla, and
S.~Narayanan, ``Linguistic analysis of differences in portrayal of movie
characters,'' in \emph{Proc.\ ACL}, 2017.

\bibitem{shafaei2020}
M.~Shafaei, N.~Safi~Samghabadi, S.~Kar, and T.~Solorio, ``Age
suitability rating: Predicting the MPAA rating based on movie
dialogues,'' in \emph{Proc.\ LREC}, 2020, pp.~1327--1335.

\bibitem{mohamed2020}
E.~Mohamed and L.~A.~Ha, ``A first dataset for film age appropriateness
investigation,'' in \emph{Proc.\ LREC}, 2020, pp.~1311--1317.

\bibitem{zhang2021}
Y.~Zhang, M.~Shafaei, F.~Gonz{\'a}lez, and T.~Solorio, ``From none to
severe: Predicting severity in movie scripts,'' in \emph{Findings of
EMNLP}, 2021, pp.~3951--3956.

\bibitem{kar2018}
S.~Kar, S.~Maharjan, and T.~Solorio, ``Folksonomication: Predicting tags
for movies from plot synopses using emotion flow encoded neural
network,'' in \emph{Proc.\ COLING}, 2018.

\bibitem{chiu2020}
M.-C.~Chiu, T.~Feng, X.~Ren, and S.~Narayanan, ``Screenplay quality
assessment: Can we predict who gets nominated?'' in \emph{Proc.\ 1st
Joint Workshop Narrative Understanding, Storylines, and Events (NUSE)},
2020.

\bibitem{sharda2006}
R.~Sharda and D.~Delen, ``Predicting box-office success of motion
pictures with neural networks,'' \emph{Expert Systems with
Applications}, vol.~30, no.~2, 2006.

\bibitem{asur2010}
S.~Asur and B.~A.~Huberman, ``Predicting the future with social media,''
in \emph{Proc.\ IEEE/WIC/ACM Web Intelligence}, 2010.

\bibitem{oghina2012}
A.~Oghina, M.~Breuss, M.~Tsagkias, and M.~de~Rijke, ``Predicting IMDb
movie ratings using social media,'' in \emph{Proc.\ ECIR}, 2012.

\bibitem{mishne2006}
G.~Mishne and N.~Glance, ``Predicting movie sales from blogger
sentiment,'' in \emph{AAAI Spring Symposium}, 2006.

\bibitem{mestyan2013}
M.~Mesty{\'a}n, T.~Yasseri, and J.~Kert{\'e}sz, ``Early prediction of
movie box office success based on Wikipedia activity big data,''
\emph{PLoS ONE}, vol.~8, no.~8, 2013.

\bibitem{madongo2024}
C.~T.~Madongo, Z.~Tang, and J.~Hassan, ``Movie box-office revenue
prediction model by mining deep features from trailers using recurrent
neural networks,'' \emph{J.\ Adv.\ Inf.\ Technol.}, 2024.

\bibitem{naeem2022}
M.~Z.~Naeem, F.~Rustam, A.~Mehmood, I.~Ashraf, and G.~S.~Choi,
``Classification of movie reviews using term frequency-inverse document
frequency and optimized machine learning algorithms,'' \emph{PeerJ
Comput.\ Sci.}, 2022.

\bibitem{tripathi2023}
J.~Tripathi, S.~Tiwari, A.~Saini, and S.~Kumari, ``Prediction of movie
success based on machine learning and Twitter sentiment analysis using
internet movie database data,'' \emph{Indonesian J.\ Elect.\ Eng.\
Comput.\ Sci.}, vol.~29, no.~3, pp.~1750--1757, 2023.

\bibitem{lee2018}
K.~Lee, J.~Park, I.~Kim, and Y.~Choi, ``Predicting movie success with
machine learning techniques: Ways to improve accuracy,'' \emph{Inf.\
Syst.\ Front.}, vol.~20, pp.~577--588, 2018.

\bibitem{bristi2019}
W.~R.~Bristi, Z.~Zaman, and N.~Sultana, ``Predicting IMDb rating of
movies by machine learning techniques,'' in \emph{Proc.\ ICCCNT}, 2019.

\bibitem{gomes2022}
A.~L.~Gomes, G.~Vianna, T.~Escovedo, and M.~Kalinowski, ``Predicting
IMDb rating of TV series with deep learning: The case of Arrow,'' in
\emph{Proc.\ Brazilian Symp.\ Inf.\ Syst.\ (SBSI)}, 2022.

\bibitem{udandarao2024}
V.~Udandarao and P.~Gupta, ``Movie revenue prediction using machine
learning models,'' \emph{arXiv:2405.11651}, 2024.

\bibitem{xie2025}
W.~Xie, ``Predicting movie success with multi-task learning: A hybrid
framework combining GPT-based sentiment analysis and SIR propagation,''
\emph{arXiv:2509.02809}, 2025.

\bibitem{vogel2020}
H.~L.~Vogel, \emph{Entertainment Industry Economics}, 10th ed.
Cambridge, U.K.: Cambridge Univ.\ Press, 2020.

\bibitem{reimers2019}
N.~Reimers and I.~Gurevych, ``Sentence-BERT: Sentence embeddings using
Siamese BERT-networks,'' in \emph{Proc.\ EMNLP}, 2019.

\bibitem{chalkidis2022}
I.~Chalkidis, X.~Dai, M.~Fergadiotis, P.~Malakasiotis, and D.~Elliott,
``An exploration of hierarchical attention transformers for efficient
long document classification,'' \emph{arXiv:2210.05529}, 2022.

\bibitem{beltagy2020}
I.~Beltagy, M.~E.~Peters, and A.~Cohan, ``Longformer: The long-document
transformer,'' \emph{arXiv:2004.05150}, 2020.

\bibitem{gorman2019}
K.~Gorman and S.~Bedrick, ``We need to talk about standard splits,'' in
\emph{Proc.\ ACL}, 2019.

\bibitem{sogaard2021}
A.~S{\o}gaard, S.~Ebert, J.~Bastings, and K.~Filippova, ``We need to
talk about random splits,'' in \emph{Proc.\ EACL}, 2021.

\bibitem{lewis2021}
P.~Lewis, P.~Stenetorp, and S.~Riedel, ``Question and answer test-train
overlap in open-domain question answering datasets,'' in \emph{Proc.\
EACL}, 2021.

\bibitem{elangovan2021}
A.~Elangovan, J.~He, and K.~Verspoor, ``Memorization vs.\
generalization: Quantifying data leakage in NLP performance
evaluation,'' in \emph{Proc.\ EACL}, 2021.

\bibitem{kapoor2023}
S.~Kapoor and A.~Narayanan, ``Leakage and the reproducibility crisis in
machine-learning-based science,'' \emph{Patterns}, vol.~4, no.~9, 2023.

\bibitem{lazaridou2021}
A.~Lazaridou \emph{et al.}, ``Mind the gap: Assessing temporal
generalization in neural language models,'' in \emph{Proc.\ NeurIPS},
2021.

\bibitem{bender2018}
E.~M.~Bender and B.~Friedman, ``Data statements for natural language
processing,'' \emph{TACL}, vol.~6, 2018.

\bibitem{dodge2021}
J.~Dodge \emph{et al.}, ``Documenting large webtext corpora: A case
study on the Colossal Clean Crawled Corpus,'' in \emph{Proc.\ EMNLP},
2021.

\end{thebibliography}

\end{document}
